Bộ dữ liệu Forest CoverType được xây dựng nhằm dự đoán loại thảm thực vật rừng
dựa trên các biến bản đồ học, không sử dụng dữ liệu viễn thám. Mỗi quan sát
tương ứng với một ô rừng kích thước 30×30 mét vuông, trong đó nhãn được xác
định từ dữ liệu của Cơ quan Lâm nghiệp Hoa Kỳ (USFS), còn các đặc trưng được
tổng hợp từ dữ liệu của USGS và USFS \cite{covertype_31}.

Khu vực nghiên cứu bao gồm bốn vùng hoang dã trong rừng quốc gia Roosevelt ở
phía bắc bang Colorado, gồm Rawah (khu vực 1), Neota (khu vực 2), Comanche Peak
(khu vực 3) và Cache la Poudre (khu vực 4). Đây là các khu vực ít chịu tác động
của con người, do đó đặc trưng thảm thực vật chủ yếu phản ánh các quá trình
sinh thái tự nhiên \cite{covertype_31}. Về đặc điểm địa hình, Neota có độ cao
trung bình lớn nhất, tiếp theo là Rawah và Comanche Peak, trong khi Cache la
Poudre có độ cao trung bình thấp nhất \cite{covertype_31}.

Xét về thành phần loài cây, Neota chủ yếu gồm spruce/fir (lớp 1); Rawah và
Comanche Peak có lodgepole pine (lớp 2) là loài chính, kèm theo spruce/fir và
aspen (lớp 5); trong khi Cache la Poudre đặc trưng bởi Ponderosa pine (lớp 3),
Douglas-fir (lớp 6) và cottonwood/willow (lớp 4) \cite{covertype_31}. Nhìn
chung, Rawah và Comanche Peak mang tính đại diện cao hơn cho toàn bộ tập dữ
liệu, còn Cache la Poudre có xu hướng khác biệt hơn do phạm vi độ cao thấp và
thành phần loài đặc trưng \cite{covertype_31}.
